\documentclass[11pt]{article}
\usepackage[a4paper,margin=25mm]{geometry}
\usepackage{amsmath,amssymb,amsthm}
\usepackage[hidelinks]{hyperref}
\emergencystretch=2em
\newtheorem{theorem}{Theorem}
\newtheorem{lemma}{Lemma}
\newcommand{\R}{\mathbb R}
\newcommand{\Z}{\mathbb Z}
\newcommand{\C}{\mathbb C}
\title{Metric lengths change the Jacobian lattice determinant\\
\large Disproof of Conjecture 00000008356}
\author{gaochengzhecpu}\date{}
\begin{document}\maketitle
\begin{abstract}
A triangle whose three edges have length $1/2$ has Jacobian cycle-lattice
Gram determinant $3/2$, whereas its ordinary spanning-tree count is $3$.
Thus the stated unweighted equality fails for metric graphs without
a unit-length normalization. The alternative convention of lattice
determinant as Euclidean covolume gives $\sqrt{3/2}$ and also fails.
\end{abstract}

\section*{Definitions and scope}
The source asserts that the determinant of the Jacobian lattice of a
metric graph equals its spanning-tree count. Here the latter has its
ordinary meaning: each spanning tree of the specified finite graph
is counted once. The source imposes no unit-length normalization.

For an oriented metric graph with positive edge lengths $\ell_e$, the
integral cycle lattice is
\[
 L=H_1(G,\Z)=\ker(\partial:\Z^{E}\longrightarrow\Z^{V}),
 \qquad Q(z,w)=\sum_{e\in E}\ell_e z_e w_e.
\]
The real cycle space modulo this lattice, with this length pairing,
is the metric Jacobian model. For an integral basis $b_1,\ldots,b_g$,
the Gram determinant is $\det(Q(b_i,b_j))$. It is independent of the
integral basis because a basis change has determinant $\pm1$.
This is the usual metric length pairing; see Mikhalkin--Zharkov,
Section 6.1 and Lemma 6.1\footnote{\href{https://arxiv.org/abs/math/0612267}{G. Mikhalkin and I. Zharkov, \emph{Tropical curves, their Jacobians and Theta functions}.}}.

\section*{The counterexample}
Let $G=K_3$ on vertices $0,1,2$, with oriented edges
$e_0:0\to1$, $e_1:1\to2$, and $e_2:2\to0$. Set every
$\ell_e=1/2$. In these coordinates,
\[
 \partial(z_0,z_1,z_2)=(z_2-z_0,z_0-z_1,z_1-z_2).
\]
Consequently $\partial z=0$ exactly when $z_0=z_1=z_2$, so
$L=\Z(1,1,1)$ and $b=(1,1,1)$ is an integral basis.
Its length Gram matrix and determinant are
\[
 [Q(b,b)]=[\tfrac12+\tfrac12+\tfrac12]=[\tfrac32],
 \qquad \det_J(G)=\tfrac32.
\]
Every spanning tree of the triangle has exactly two edges, and every
two-edge subset is a spanning tree. There are exactly three such
subsets. Hence
\[
 \boxed{\det_J(G)=\tfrac32\ne3=\tau(G).}
\]
If determinant denotes Euclidean covolume instead, the left side is
$\sqrt{3/2}\ne3$, since equality would imply $3/2=9$.

The missing lengths are essential. For this triangle the metric
tree-complement expression is
$\sum_T\prod_{e\notin T}\ell_e=3/2$; it counts each omitted edge
with its length and is not the unweighted tree count $3$.
With unit edge lengths the Gram determinant would indeed be $3$.
Thus the counterexample concerns exactly the unnormalized metric-graph
equality stated in the source. Replacing its right side by a weighted
enumerator gives a different statement. Failure of the first asserted
equality suffices; the other clauses are not needed.

\section*{Correspondence with Lean}
The Lean object \texttt{K3} is an actual \texttt{SimpleGraph}.
The subtype \texttt{SpanningTree} consists of its spanning connected
acyclic subgraphs, and \texttt{spanning\_\allowbreak{}tree\_\allowbreak{}count} proves its
cardinality is $3$ via an exhaustive finite edge-mask bijection.
The theorem \texttt{oriented\_\allowbreak{}edges\_\allowbreak{}are\_\allowbreak{}actual} checks that the
chosen oriented edges enumerate the edges of that same graph.

The linear map \texttt{boundary} comes from the actual incidence
matrix. Its kernel \texttt{CycleLattice} is equipped with an explicit
linear equivalence to $\Z$, giving the actual Mathlib
\texttt{Basis} named \texttt{cycleBasis}. The matrix \texttt{gram}
evaluates the edge-length pairing on this basis.
The final theorems \texttt{determinant\_\allowbreak{}ne\_\allowbreak{}spanning\_\allowbreak{}tree\_\allowbreak{}count}
and \texttt{covolume\_\allowbreak{}ne\_\allowbreak{}spanning\_\allowbreak{}tree\_\allowbreak{}count} establish both
convention-specific inequalities. The formalized object is the
metric cycle lattice; no unrelated scalar is assigned the name
``Jacobian determinant.''

\section*{Reproduction}
The project pins Lean 4.19.0 and Mathlib commit
\\\texttt{c44e0c8ee63ca166450922a373c7409c5d26b00b}.\par
From \texttt{lean/}, run \texttt{lake build} and
\texttt{lake env lean -DwarningAsError=true Main.lean}.
On a new machine, \texttt{lake exe cache get} retrieves the official
pinned dependency artifacts. The submission itself is checked in a fresh
build directory. Printed dependencies use only the standard
\texttt{propext}, \texttt{Classical.choice}, and \texttt{Quot.sound}.
No additional axioms, proof placeholders, or numerical oracles are used.
Run \texttt{tectonic main.tex} to reproduce the PDF; no auxiliary numerical
script is needed. Build logs, source hashes, and visual-review records
are included under \texttt{verification/}.
\section*{Conjecture source}
\href{https://github.com/The-Last-Math-Competition/The-Last-Math-Competition/blob/main/conjectures/00000008356.md}{The Last Math Competition, Conjecture 00000008356}. The exact bilingual source and its commit and SHA-256 are preserved in this submission.
\end{document}
